\chapter{Japanese Encephalitis}

\section*{Definition and Overview}
Japanese encephalitis (JE) is a viral infection of the brain caused by the Japanese encephalitis virus, which is spread to humans by the bite of infected mosquitoes, mainly of the Culex species. It occurs across much of Asia and the western Pacific, and was first reported in mainland Australia in 2022, when cases occurred in south-eastern Australia. Most people infected with the virus have no symptoms or a mild flu-like illness, but in a minority the virus infects the brain, causing high fever, headache, confusion, seizures, and paralysis, with death or lasting disability in a substantial proportion of severe cases. There is no specific antiviral treatment, so prevention through vaccination and avoiding mosquito bites is essential.

\section*{Epidemiology}
Japanese encephalitis is the most important cause of viral encephalitis in Asia, causing an estimated 70,000 cases and 10,000 to 15,000 deaths each year, with the highest rates in rural and agricultural areas of China, India, South-East Asia, and the western Pacific. It mainly affects children in endemic areas. The virus is transmitted between pigs and wading birds by Culex mosquitoes, and humans are infected when they are bitten in or near these transmission cycles. In Australia, cases had previously only occurred in travellers, but locally acquired cases appeared in 2022 in New South Wales, Victoria, South Australia, and Queensland, and a vaccine has now been funded for at-risk groups.

\section*{Aetiology and Pathophysiology}
The Japanese encephalitis virus is a flavivirus, related to dengue, Zika, and West Nile viruses. It is transmitted when a Culex mosquito, usually Culex annulirostris in Australia, bites an infected animal, most often pigs or waterbirds, and then bites a human. The virus multiplies in the skin and lymph nodes and spreads through the blood to the brain, where it infects nerve cells and triggers inflammation. Most infections are contained by the immune system, but when the virus crosses the blood-brain barrier it causes encephalitis with swelling, nerve cell damage, and, in severe cases, destruction of brain tissue. The incubation period is 5 to 15 days.

\section*{Clinical Features}
\begin{itemize}
    \item Most infections cause no symptoms or a mild illness with fever and headache
    \item Encephalitis: high fever, severe headache, neck stiffness, and vomiting
    \item Confusion, drowsiness, and altered consciousness
    \item Seizures, which are common in children
    \item Tremor, muscle weakness, and paralysis
    \item Stiff neck and sensitivity to light
    \item Coma in severe cases
    \item Neurological deficits that may persist after recovery, including movement and cognitive problems
\end{itemize}

\section*{Differential Diagnosis}
Other causes of encephalitis and similar illnesses must be considered.
\begin{itemize}
    \item \textbf{Other mosquito-borne viruses:} Murray Valley encephalitis, West Nile, and dengue
    \item \textbf{Herpes simplex encephalitis:} a treatable viral encephalitis
    \item \textbf{Meningitis:} bacterial or viral infection of the meninges
    \item \textbf{Other flaviviruses:} Kunjin and other Australian viruses
    \item \textbf{Non-infectious causes:} stroke, brain tumours, and autoimmune encephalitis
    \item \textbf{Severe systemic infections:} malaria and sepsis with brain involvement
\end{itemize}

\section*{When to Seek Medical Care}
Anyone with fever and headache after travel to or residence in a risk area should be assessed, and fever with confusion, seizures, or neurological symptoms is an emergency. People travelling to JE-endemic regions should seek advice about vaccination before departure, and Australians living or working in affected pig or mosquito areas should consider vaccination.

\textbf{Call 000 (or local emergency services) for:}
\begin{itemize}
    \item High fever with confusion, drowsiness, or a seizure
    \item Stiff neck, severe headache, and sensitivity to light
    \item Weakness, paralysis, or difficulty speaking
    \item Collapse or loss of consciousness
    \item Difficulty breathing
\end{itemize}

\section*{Diagnosis and Investigations}
\begin{itemize}
    \item \textbf{Blood tests:} antibodies to the virus, with repeat testing to confirm recent infection
    \item \textbf{Cerebrospinal fluid analysis:} lumbar puncture to examine the fluid around the brain
    \item \textbf{PCR testing:} detection of the virus in blood or cerebrospinal fluid early in the illness
    \item \textbf{MRI of the brain:} characteristic changes in the thalami and other regions
    \item \textbf{Electroencephalogram (EEG):} in people with seizures
    \item \textbf{Exclusion of other causes:} testing for other viruses, bacteria, and non-infectious causes
\end{itemize}

\section*{Management}
\begin{itemize}
    \item \textbf{No specific antiviral treatment:} care is supportive
    \item \textbf{Hospital care:} most people with encephalitis require hospitalisation
    \item \textbf{Seizure management:} anticonvulsant medicines
    \item \textbf{Brain swelling management:} medicines and monitoring to reduce pressure
    \item \textbf{Respiratory and intensive care:} breathing support in severe cases
    \item \textbf{Rehabilitation:} physiotherapy, speech therapy, and cognitive support after recovery
    \item \textbf{Prevention for others:} mosquito control and vaccination programs
    \item \textbf{Public health follow-up:} notification and investigation of cases
\end{itemize}

\section*{Complications}
\begin{itemize}
    \item Death in around 20 to 30 per cent of people with clinical encephalitis
    \item Permanent neurological disability: paralysis, movement disorders, cognitive impairment, and seizures
    \item Behavioural and psychological problems after recovery
    \item Prolonged hospitalisation and intensive care
    \item Effects on families and caregivers
\end{itemize}

\section*{Prognosis and Outcomes}
The outcome depends on the severity of the infection. Around one in four people with clinical encephalitis dies, and up to half of survivors have lasting neurological problems. Among people with mild or silent infection, there are no long-term effects. Because there is no specific treatment, prevention is the key: vaccination is highly effective, and avoiding mosquito bites reduces exposure. In Australia, funded vaccination is available for people at increased risk, and surveillance and mosquito control programs continue.

\section*{Prevention and Self-Care}
\begin{itemize}
    \item Get vaccinated against Japanese encephalitis if you live in, travel to, or work in at-risk areas
    \item Use insect repellent containing DEET or picaridin
    \item Wear long sleeves, long pants, and covered shoes outdoors at dawn and dusk
    \item Sleep under mosquito nets and use screens and insecticide where needed
    \item Remove standing water around the home where mosquitoes breed
    \item Check travel health advice before going to Asia or the western Pacific
    \item Report fever and headache after travel or exposure promptly
    \item Support mosquito surveillance and control programs in affected areas
\end{itemize}

\section*{Sources and Further Reading}
The following reputable sources provide further information for healthcare students and patients seeking reliable, current advice about Japanese encephalitis.
\begin{itemize}
    \item Department of Health and Aged Care. \textit{Japanese encephalitis information and vaccine advice.} https://www.health.gov.au/
    \item Australian Immunisation Handbook. \textit{Japanese encephalitis vaccine.} https://immunisationhandbook.health.gov.au/
    \item Healthdirect Australia. \textit{Japanese encephalitis.} https://www.healthdirect.gov.au/
    \item NSW Health. \textit{Japanese encephalitis fact sheet.} https://www.health.nsw.gov.au/
    \item World Health Organization. \textit{Japanese encephalitis.} https://www.who.int/
    \item Centers for Disease Control and Prevention. \textit{Japanese encephalitis.} https://www.cdc.gov/
\end{itemize}
